% bank_fintech_substance/drafts/en/manuscript.tex
% P3 manuscript skeleton — M0 G0 scaffolded per B-route plan v2.55 (2026-05-17)
% Compile: xelatex + biber + xelatex × 2

% papers_assets/_shared/preamble_en.tex
% Shared English LaTeX preamble for all 5 paper drafts.
% P4.1 (plan-r3 derivative-paper extension, 2026-05-15)
% Compile: xelatex + biber + xelatex × 2

\documentclass[12pt,a4paper]{article}

% --- Fonts & Encoding ---
\usepackage{fontspec}
\setmainfont{Times New Roman}[Ligatures=TeX]

% --- Math ---
\usepackage{amsmath, amssymb, amsthm}

% --- Tables & Figures ---
\usepackage{booktabs}
\usepackage{graphicx}
\usepackage{longtable}
\usepackage{rotating}
\usepackage{caption}
\captionsetup{font=small,labelfont=bf,skip=4pt}

% --- Layout ---
\usepackage[margin=1in]{geometry}
\usepackage{setspace}
\onehalfspacing
\usepackage{hyperref}
\hypersetup{colorlinks=true,linkcolor=blue,citecolor=blue,urlcolor=blue}

% --- Bibliography (biblatex + biber) ---
\usepackage[backend=biber,style=apa,natbib=true]{biblatex}

% --- Useful macros ---
\newcommand{\code}[1]{\texttt{#1}}
\newcommand{\NoteThis}[1]{\textbf{\textcolor{red}{[TODO: #1]}}}
\usepackage[dvipsnames]{xcolor}

% --- Theorem environments ---
\newtheorem{prop}{Proposition}
\newtheorem{hyp}{Hypothesis}


\addbibresource{refs.bib}

\title{Substantive vs Symbolic Fintech Strategies:\\Long-Run ROA Reversal Evidence from Chinese Banking}
\author{\NoteThis{Author names}}
\date{\today}

\begin{document}
\maketitle

\begin{abstract}
\NoteThis{Final abstract pending — target $\approx$ 200 English words.}
This paper applies Westphal--Zajac's (1994, 1998 \emph{ASQ}) substantive vs symbolic
management framework to bank fintech strategic announcements during 2007--2024. We
classify each bank-year fintech announcement using a binary substantiveness measure
(\code{innovation\_substantiveness\_full}) derived from large-language-model annotation,
and test the dual-outcome prediction: symbolic announcements should produce positive
short-run market reactions but long-run ROA reversal (within-bank 3-year forward change),
whereas substantive announcements should not. Identification leverages four exogenous
policy shocks (余额宝 launch, P2P regulation, asset-management reform, Ant IPO pause) for
cross-sectional intensity DiD. We find \NoteThis{key result}\@.
\end{abstract}

\section{Introduction}\label{sec:intro}
\NoteThis{Motivate the research question.}

This study makes three contributions. First, we operationalize Westphal--Zajac's
substantive/symbolic dichotomy in the fintech-strategy context using large-scale media
text and LLM annotation, complementing prior survey-based and case-based evidence.
Second, we exploit an 18-year Chinese banking panel that contains complete short-run
event windows and long-run ROA evolution, allowing dual-outcome testing of the W\&Z
implicit mechanism. Third, we differentiate measurement constructs from concurrent
fintech-tone and fintech-attention studies, with explicit anti-salami-slicing audit
(per replication package).

\section{Literature Review}\label{sec:lit}
\NoteThis{Three strands: (1) Westphal--Zajac substantive/symbolic
\autocite{westphal1994substance,westphal1998sym,westphal2001decoupling};
(2) Symbolic management of strategic change \autocite{fisszajac2006symbolic};
(3) Practice-based view of strategy \autocite{bromiley2014practice}.}

\section{Methodology}\label{sec:method}
\subsection{Sample and Data}
Bank--year panel covers $N \times T \approx 3000\text{--}9000$ bank--years across
2007--2024 (per spec.sample). All bank--year observations are processed through the
\code{news\_process} pipeline (stage 32a innovation\_substantiveness).

\subsection{Substantiveness Measurement}
\NoteThis{Describe LLM annotation pipeline producing innovation\_substantiveness\_full
binary indicator (W\&Z substantive vs symbolic); cross-validate with full vs ctx tracks.}

\subsection{Dual-Outcome Identification}
\NoteThis{Long\_run\_roa\_change\_3yr (within-bank 3-year forward) + car\_3day\_window
\_fintech\_announcement (short-run); within-bank within-year FE; cluster at bank level.}

\section{Data}\label{sec:data}
\NoteThis{Data sources, sample selection, descriptive statistics (Table 1).}

\section{Results}\label{sec:results}
\NoteThis{Table 2 (M1$\to$M6 progressive main regression); Table 3 (dual outcome:
short-run CAR + long-run ROA reversal jointly); Table 4 (heterogeneity 4 bank types
$\times$ SOE $\times$ fintech wave $\times$ capital\_ratio quartile); Table 5
(robustness: alternative substantiveness track full vs ctx).}

\section{Conclusion}\label{sec:conclusion}
\NoteThis{Summarize findings and implications for strategic management theory.}

\printbibliography

\end{document}
